% ConTeXt document produced by My Document Converter
\setuppapersize[A4]
\setuplayout[width=middle,height=middle,backspace=20mm,topspace=20mm]
\setupbodyfont[11pt]
\setupwhitespace[medium]
\setuphead[section,subsection,subsubsection][style=\tfa\bf]
\setupinteraction[state=start,color=blue]
\starttext
\startalignment[center]
\blank[2*big]
{\tfd My Document Converter 샘플 문서}
\blank[3*medium]
{\tfa SHKWON}
\blank[2*medium]
{\tfa 2026-09-22}
\blank[3*medium]
\stopalignment

\startsubject[title={소개 (Introduction)},reference={소개-introduction}]

이 문서는 {\bf My Document Converter} 가 지원하는 모든 출력 형식의 샘플을 만드는 데 쓰이는 원본입니다.
{\em 강조}, {\bf 굵게}, \overstrikes{취소선}, \type{인라인 코드}, H\low{2}O 와
x\high{2}, 그리고 \useURL[url1][https://example.com][][링크]\from[url1]를 담고
있습니다.
This paragraph is in English so that Latin-only formats still show
something readable.

\startsubsubject[title={목록 (Lists)},reference={목록-lists}]

\startitemize[packed]
\item 첫 번째 항목
\item 두 번째 항목
\startitemize[packed]
\item 중첩된 항목
\item 또 하나의 중첩 항목
\stopitemize
\item 세 번째 항목
\stopitemize

\startitemize[n,packed]
\item 순서가 있는 항목
\item 두 번째
\item 세 번째
\stopitemize

\startitemize[packed]
\item \boxplus 끝난 일
\item \square 남은 일
\stopitemize

\startdescription{용어}
용어에 대한 정의입니다.
\stopdescription
\startdescription{Pandoc}
범용 문서 변환기. 이 프로그램은 같은 형식 집합을 자체 엔진으로 다룹니다.
\stopdescription

\stopsubsubject

\startsubsubject[title={코드 (Code)},reference={코드-code}]

\starttyping[option=js]
function greet(name) {
  console.log(`Hello, ${name}!`)
}
greet('world')
\stoptyping

\stopsubsubject

\startsubsubject[title={표 (Table)},reference={표-table}]

\startplacetable[title={지원 형식의 일부}]
\startxtable
\startxtablehead[head]
\startxrow
\startxcell[align=right] {\bf 형식} \stopxcell
\startxcell[align=middle] {\bf 확장자} \stopxcell
\startxcell[align=middle] {\bf 읽기} \stopxcell
\startxcell[align=left] {\bf 쓰기} \stopxcell
\stopxrow
\stopxtablehead
\startxtablebody[body]
\startxrow
\startxcell[align=right] Markdown \stopxcell
\startxcell[align=middle] .md \stopxcell
\startxcell[align=middle] O \stopxcell
\startxcell[align=left] O \stopxcell
\stopxrow
\startxrow
\startxcell[align=right] HTML \stopxcell
\startxcell[align=middle] .html \stopxcell
\startxcell[align=middle] O \stopxcell
\startxcell[align=left] O \stopxcell
\stopxrow
\startxrow
\startxcell[align=right] DOCX \stopxcell
\startxcell[align=middle] .docx \stopxcell
\startxcell[align=middle] O \stopxcell
\startxcell[align=left] O \stopxcell
\stopxrow
\startxrow
\startxcell[align=right] EPUB \stopxcell
\startxcell[align=middle] .epub \stopxcell
\startxcell[align=middle] O \stopxcell
\startxcell[align=left] O \stopxcell
\stopxrow
\stopxtablebody
\stopxtable
\stopplacetable

\stopsubsubject

\startsubsubject[title={인용과 수식 (Quote and math)},reference={인용과-수식-quote-and-math}]

\startblockquote
좋은 문서는 어떤 형식으로도 읽힙니다.
— 누군가
\stopblockquote

인라인 수식 $E = mc^2$ 과 블록 수식:

\startformula
\int_0^1 x^2\,dx = \frac{1}{3}
\stopformula

\stopsubsubject

\startsubsubject[title={각주와 그림 (Footnote and figure)},reference={각주와-그림-footnote-and-figure}]

각주가 붙은 문장입니다.\footnote{각주의 내용입니다.} 자동 링크
\useURL[url2][https://pandoc.org][][https://pandoc.org]\from[url2] 와
\useURL[url3][mailto:knix008@naver.com][][mailto:knix008@naver.com]\from[url3]
도 있습니다.

\externalfigure[../public/app-icon.svg]

\thinrule

마지막 문단입니다. 줄 끝에 두 칸을 두면\crlf
강제 줄바꿈이 됩니다.

\stopsubsubject

\stopsubject

\stoptext
